\section{Implementation}\label{sec:impl}

Bellwether is a Rust workspace with one crate per work package
(\doc{work-packages.md}), each with one owner and one interface contract. The
contract crate,
\code{bellwether-types}, defines every record, trait, and format constant and
contains no logic beyond construction, hashing, and serialization; changing it
is an architectural decision. Table~\ref{tab:crates} lists the crates and their
status at the time of writing. Each crate lands by a pull request reviewed
against the invariants and decision records it owns.

\begin{table}[t]
  \caption{Crates and their status at the time of writing. \emph{Merged}: on
    the main branch. \emph{Review}: an open pull request. \emph{Planned}:
    specified, not yet written.}
  \label{tab:crates}
  \small
  \begin{tabular}{@{}l>{\raggedright\arraybackslash}p{0.53\columnwidth}l@{}}
    \toprule
    Crate & Role & Status \\
    \midrule
    types & records, traits, format constants & merged \\
    source & Parquet and chunked text sources; fingerprints & merged \\
    ledger & content-addressed records; write-time rules & merged \\
    core & canonical encoding, segments, bloom filters, bitmaps, temporal
      chains, merge & review \\
    store & conditional PUT, manifest commits, cache & review \\
    bench & synthetic ledger, workloads, harness, SLOs & review \\
    build & discovery, id assignment, segment build, compaction & planned \\
    profile & fanout, sortedness, lints, candidate composites & planned \\
    replay & workload recorder; exact replay & planned \\
    verify & reconciliation; Certificates & planned \\
    cli, mcp & one shared code path, two interfaces & planned \\
    evolve & the optimizer of Section~\ref{sec:evolve} & planned \\
    \bottomrule
  \end{tabular}
\end{table}

\paragraph{Formats.} Storage, record, and encoding formats are versioned from
the first commit, and each is at version 1 (\doc{format.md}). A binary that
meets a newer version refuses with a typed error and never downgrades in place
(\adr{017}). The canonical encoding gives each Arrow type one byte form:
integers widened to 64 bits big-endian, with the sign bit flipped for signed
types, floats under a total order with one canonical NaN, decimals rescaled to the schema's declared scale
(an unrepresentable value is an error, never rounded), escaped and terminated
strings, and timestamps normalized to UTC in the schema's unit (\adr{016}).

\paragraph{Memory bounds and budgets.} Memory is bounded per operation
(\adr{018}): a build streams one row group per worker and holds at most one
segment, whose name is the hash of its bytes; a lookup holds one block index and
the bitmaps for the keys asked; a replay holds one run's key set. Observe tools
have a 100\,ms budget (\doc{architecture.md}). For evolve on the medium
synthetic dataset, the budgets are 5\,ms per replay of one run under one
candidate, one second to attribute a 20-node plan, one minute per generation of
16 candidates excluding proposer latency, and one hour for 50 generations
(\doc{evolve.md} \S12). The harness asserts these budgets; none is a
measurement.

\paragraph{A preliminary microbenchmark.} The core crate's review branch
includes a Criterion benchmark of single-key lookup in an exact segment of
64-bit integer keys held in memory. It reports about 0.2\,\textmu s per lookup
with $10^4$ keys and about 0.5\,\textmu s with $10^6$ keys. The caveats are
material: the crate is not merged; the segment is in memory, so storage I/O,
segment open, and the manifest read are excluded; the key is a single column; no
machine description accompanies the number, although the benchmark methodology
requires one for every reported result (\doc{benchmark.md}); and in-memory
lookup latency is not one of the paper's metrics. We report the number as
context and draw no conclusion from it.

\paragraph{Interfaces.} Command-line verbs such as \code{why}, \code{fix},
\code{latest}, \code{asof}, and \code{reconcile}, with \code{evolve} and
\code{report} for the optimizer, share one code path with the MCP tiers of
Section~\ref{sec:loop} (\doc{work-packages.md}).
